\documentclass{article}

\textwidth=6.5in
\textheight=8.5in
\voffset=-54pt
\oddsidemargin=0in
\evensidemargin=0in
\setlength{\parskip}{8pt}
\setlength{\parindent}{0pt}

\usepackage{amsmath, amssymb}
\usepackage{ifthen}

\usepackage[inline]{enumitem}
\setlist{topsep=0pt, parsep=8pt}

\usepackage[rgb,table]{xcolor}\convertcolorsUfalse
\usepackage{graphicx}

% change hyperref options
\PassOptionsToPackage{hyphens}{url}
\usepackage[bookmarks,colorlinks,linkcolor=black,citecolor=blue,pdfstartview=FitH,urlcolor=blue]{hyperref}
\hypersetup{pdfpagemode=UseNone}
\usepackage{bookmark}

\usepackage{graphicx}
\usepackage{multicol}
\usepackage{framed}
\usepackage{array}

\usepackage{icomma}
\usepackage{pifont}

\usepackage{marginnote}
\usepackage[colorinlistoftodos]{todonotes}
\renewcommand{\marginpar}{\marginnote}

\usepackage{soul}

\usepackage{pgfplots}
\pgfplotsset{compat=1.18}  
\pgfplotscreateplotcyclelist{mycolor}{black, blue, red, green!70!black, magenta, purple, cyan, orange }  
\pgfplotsset{
  disabledatascaling,
  cycle list=tikzcolor, 
  axis lines=middle,
  every axis plot/.style={no marks, thick},
  every axis/.style={
    enlarge x limits=false,
    enlarge y limits=auto,
    xlabel style={at={(current axis.right of origin)},anchor=west},
    ylabel style={at={(current axis.above origin)},anchor=south},
    % extra x and y ticks for asymptotes
    extra x tick style={grid=major, grid style={dashed,black}},
    extra y tick style={grid=major, grid style={dashed,black}},
    /pgf/number format/1000 sep={},
  },    
  tick label style = {font=\footnotesize, fill=\currentbgcolor, inner sep=0pt},    
  xticklabel shift=1pt,    
  scaled ticks=false,
  scaled y ticks=false,
  unbounded coords=jump,
  samples=101,
  minor tick length=0,
  scatter/classes={
    open={mark=*,draw=black,fill=white, mark size=2, thin},
    closed={mark=*,draw=black,fill=black, mark size=2, thin}
  },
  width=3in,
}
\pgfplotsset{open/.style={only marks, mark=*, fill=white, thin, mark size=2}}
\pgfplotsset{closed/.style={only marks, mark=*, draw=black, fill=black,
    thin, mark size=2}}
\pgfplotsset{labeled/.style={nodes near coords, only marks, mark=*, draw=black, fill=black, thin, mark size=2, point meta=explicit symbolic}}

\usepackage{tikz}
\usetikzlibrary{
  matrix,%
  % enables the snake paths
  decorations.pathmorphing,%
  % enables bigger arrows
  decorations.markings,%
  decorations.pathreplacing,%
  decorations.text,%
  calc,%
  patterns,%
  shapes.geometric,%
  arrows,
  decorations.shapes,
  positioning,
  plotmarks,
  intersections
}
\tikzset{>=latex} % sets default arrow type in tikz
\tikzset{font=\normalsize}
\tikzset{mark size=1.5pt, mark options=thin}
\tikzset{pin distance=4pt,
  every pin edge/.style={<-, thin, shorten <= -2pt}}

\interfootnotelinepenalty=10000
\renewcommand{\thefootnote}{(\arabic{footnote})}

\usepackage{multirow, bigdelim}

% change to \usepackage[sols]{handout} to see solutions
\usepackage[sols, grading, workshop]{handout}

\newcommand\iscurrentenumi[1]{%
  \ifthenelse{\equal{#1}{\theenumi}}{}{#1}}
\setenumerate[1]{ref=\#\arabic*}
\setenumerate[2]{ref=\protect\iscurrentenumi{\theenumi}(\alph*)}
\newcommand\enumiiprefix[2]{%
  \ifthenelse{{\equal{#1}{\theenumi}}\AND{\equal{#2}{(\alph{enumii})}}}{}{}%
  \ifthenelse{{\equal{#1}{\theenumi}}\AND{\NOT{\equal{#2}{(\alph{enumii})}}}}{#2}{}%
  \ifthenelse{\equal{#1}{\theenumi}}{}{#1#2}%
}
\setenumerate[3]{ref=\protect\enumiiprefix{\theenumi}{(\alph{enumii})}\roman*}

\makeatletter

% underline links               
\let\@href=\href
\renewcommand{\href}[2]{\@href{#1}{\ul{#2}}}

\makeatother

\definecolor{purple}{rgb}{0.5,0,1.0}

\newcommand{\doedfinity}[2][\newpage]{
  \begin{problem}[#1]
    Do the Edfinity assignment ``Problem Set #2''.

    We encourage you to write down any scratch work here, as well as
    things you want your future self to remember. Nothing you write
    here will be graded; it's just for you!
  \end{problem}
}

\newcommand{\psrnewpage}{\newpage}
\newcommand{\psreflection}[2][]{

  \ifthenelse{\not\equal{#1}{nocite}}{
    \begin{enumerate}[resume]
    \item
      \begin{problem}[1in]
        Please cite any help you got on this problem set:
        \begin{itemize}[nosep]
        \item If you worked with other students, please list their names here.
        \item If you used resources other than those on our Canvas site, please list them here.
        \end{itemize}
      \end{problem}

    \end{enumerate}
  }

  \probonly{%
    #2

    \psrnewpage

    \emph{Reflection space.} It's a good habit to take a few minutes at the end of each pset to reflect on what you learned and jot down notes to your future self. Here are a few reflection questions to get you started. Nothing you write here will be graded; it's just for you!
    \begin{itemize}
    \item Take a few minutes to look back at the problems. How could you explain to someone else how to do each problem? What was your strategy, and how did you know to use that strategy? If there are problems where you don't feel able to answer this, write down which ones they are. Come back to them in a few days when we post the homework solutions!

      \vspace{1.5in}

    \item Take a look back at the learning objectives on today's worksheet; how are you feeling about each one? If there are ones you know you need more practice with, write those down here.

      \vspace{1.5in}

    \item What did you find most challenging about this material? Do you have any lingering questions about it? If so, write them down here so you don't forget them. At some point, stop by office hours or MQC to ask!

      \vfill
    \end{itemize}      
  }
}

\usepackage{fancyhdr}
\lhead{\textcolor{white}{This content is protected and may not be shared, uploaded, or distributed.}}
\rhead{}
\rfoot{\vspace*{\fill}\footnotesize\textcopyright{} \emph{Harvard Math Ma}}
\renewcommand{\headrulewidth}{0pt}
\pagestyle{fancy}
\begin{document}

\begin{center}
  \large\textsc{Problem Set 32 - Solving Equations with Logarithms and Exponents}
\end{center}

\begin{enumerate}
\item  \note{
    \begin{itemize}[nosep]
    \item Match basic functions (the ones from the beginning of the semester, plus exponentials and logs) to graphs
    \item Solve the 3 basic types of equations (similar to {{\#3}} on \href{https://people.math.harvard.edu/~jjchen/mathma/docs/worksheets/worksheet32-sols.html}{the worksheet ``Solving Equations with Logarithms and Exponents''})
    \item Solve some very basic modifications of those (involving $ax + b$ instead of just $x$)
    \item Pick correct log / exp rules from a list
    \end{itemize}
}

  \doedfinity{32}  

\item \label{exp-log-solve-nice-answers}

  The following equations have nice solutions (integers or simple fractions). Solve each equation without using technology; please organize your work to make it easy for the reader to follow.

  Plug your solutions back in to the original equation to check that they really work. (Please write down your work for this.) 

  \begin{enumerate}

  \item

    \begin{grading}
      4 points:
      \begin{itemize}[nosep]
      \item 1 for correctly getting to $x = (2 - x)^2$
      \item 1.5 for correctly solving $x = (2 - x)^2$
      \item 1 for recognizing that only $x = 1$ is in the domain
      \item 0.5 for checking that $x = 1$ works
      \end{itemize}
    \end{grading}

    \begin{problem}[\vfill]
      $\ln x = 2 \ln(2 - x)$
    \end{problem}

    \begin{solution}
      \begin{align*}
        \ln x &= 2\ln(2-x) \\
        \ln x &= \ln[(2-x)^2] \\
        x &= (2-x)^2 \\
        x &= x^2 - 4x + 4 \\
        0 &= x^2 - 5x + 4 \\
        0 &= (x-1)(x-4) 
      \end{align*}
      This gives us $x=1$ and $x=4$. Plugging in $x=1$, we see 
      that it is a solution.
      \begin{align*}
        \ln 1 &= 2\ln(2-1)\\
        0 &= 0\ \text{\ding{51}}\\
        \intertext{Plugging in $x=4$, we see that it is extraneous.}
        \ln 4 &= 2\ln(2-4)\\
        \ln 4 &= 2 \ln(-2)\ \text{\ding{55}}
      \end{align*}
      So the only solution is $x=\boxed{1}$.
    \end{solution}

  \item

    \begin{grading}
      3.5 points:
      \begin{itemize}[nosep]
      \item 1.5 for getting to $x^2 = 2 - x$. (If a student starts out trying to take $\log_5$ of both sides, like $\log_5 \left( 5^{x^2} - \dfrac{25}{5^x} \right) = \log_5(0)$, they should lose 1.5 points, since they can't simplify the left side and $\log_5(0)$ is undefined.)
      \item 1 for solving the quadratic $x^2 = 2 - x$
      \item 1 for checking that both solutions work
      \end{itemize}
    \end{grading}

    \begin{problem}[\vfill\newpage]
      $5^{x^2} - \dfrac{25}{5^x} = 0$
    \end{problem}

    \begin{solution}
      We can rewrite the equation so that we have one term
      on each side and take the log base 5 of both sides.
      \begin{align*}
        5^{x^2} - \frac{25}{5^x} &= 0 \\
        5^{x^2} &= \frac{25}{5^x} \\
        \log_5(5^{x^2}) &= \log_5\!\left(\frac{25}{5^x}\right) \\
        \log_5(5^{x^2}) &= \log_5 25 - \log_5 (5^x) \\
        x^2 &= 2 - x \\
        x^2 +x - 2 &= 0 \\
        (x+2)(x-1) &= 0
      \end{align*}
      This gives us the solutions $x=-2$ and $x=1$. Plugging
      in $x=-2$, we see that it is a solution.
      \begin{align*}
        5^{(-2)^2} - \frac{25}{5^{-2}} &= 0 \\
        5^4 - 5^4 &= 0\ \text{\ding{51}} \\
        \intertext{Plugging in $x=1$, we see that it is
        also a solution.}
        5^{1^2} - \frac{25}{5^1} &= 0\\
        5 - 5 &= 0 \ \text{\ding{51}}
      \end{align*}
      So the two solutions are $x=\boxed{-2}$ and $x=\boxed{1}$.
    \end{solution}

  \end{enumerate}

\item  In this problem, you'll work through 2 ways to solve the following equation  for $x$.
  \[ \ln \left( \sqrt{x} \right) + \ln \left( ex^3 \right) - 2\ln x = 4\]

  \begin{enumerate}

  \item

    \begin{grading}
      3 points:
      \begin{itemize}[nosep]
      \item 0.5 for knowing $2 \ln x = \ln \left( x^2 \right)$
      \item 0.5 for correctly rewriting $\ln(\text{stuff}) + \ln(\text{stuff})$
      \item 0.5 for correctly rewriting $\ln(\text{stuff}) - \ln(\text{stuff})$
      \item 0.5 for correctly raising $e$ to both sides
      \item 1 for solving for $x$
      \end{itemize}
    \end{grading}

    \begin{problem}[\stretch{1.25}]
      Here's one way. First use log rules to rewrite the left hand side as a single logarithm (so the equation looks like $\ln(\text{stuff}) = 4$). Then raise $e$ to both sides of the equation, and find $x$. Please simplify your answer.
    \end{problem}

    \begin{solution}
      \begin{align*}
        \ln\!\left( \sqrt{x} \right) + \ln\!\left( ex^3 \right) - 2\ln x 
        &= 4 \\
        \ln\!\left( \sqrt{x} \right) + \ln\!\left( ex^3 \right) - 
        \ln\!\left( x^2 \right) &= 4 \\
        \ln\!\left(\frac{\sqrt{x} \cdot ex^3}{x^2}\right) &= 4 \\
        \ln\!\left(ex^{3/2}\right) &= 4 \\
        e^{\ln(ex^{3/2})} &= e^4 \\
        ex^{3/2} &= e^4 \\
        x^{3/2} &= e^3 \\
        x &= (e^3)^{2/3} \\
        x &= \boxed{e^2}
      \end{align*}
    \end{solution}

  \item

    \begin{grading}
      2.5 points:
      \begin{itemize}[nosep]
      \item 0.5 for knowing $\ln \left( \sqrt{x} \right) = \dfrac{1}{2} \ln x$
      \item 0.5 for knowing $\ln \left( ex^3 \right) = \ln e + \ln \left( x^3 \right)$
      \item 0.5 for knowing $\ln \left( x^3 \right) = 3 \ln x$
      \item 0.5 for combining terms on the left correctly to get to $\ln x = 2$ 
      \item 0.5 for knowing $x = e^2$
      \end{itemize}
    \end{grading}

    \begin{problem}[\vfill]
      Here's a second way. Use log rules to rewrite each term on the left hand side in terms of $\ln x$. Then simplify the left hand side, solve for $\ln x$, and find $x$. 
    \end{problem}

    \begin{solution}
      \begin{align*}
        \ln\!\left( \sqrt{x} \right) + \ln\!\left( ex^3 \right) - 2\ln x 
        &= 4 \\
        \ln\!\left( \sqrt{x} \right) + + \ln e + 
        \ln\!\left(x^3 \right) - 2\ln x &= 4 \\
        \tfrac{1}{2}\ln x + 1 + 3\ln x - 2 \ln x &= 4 \\
        \tfrac{3}{2}\ln x + 1 &= 4 \\
        \tfrac{3}{2}\ln x &= 3 \\
        \ln x &= 2 \\
        x &= \boxed{e^2}
      \end{align*}
    \end{solution}

  \item

    \begin{grading}
      0.5 points, any answer is fine.
    \end{grading}

    \begin{problem}[\nstretch{0.25}]
      Which way did you prefer? 
    \end{problem}

    \begin{solution}
      This is totally up to you!
    \end{solution}

  \end{enumerate}

\item % 13.3.4

  Solve for $x$ without using technology; give exact answers. Please organize your work so that your reader can easily tell what you did in each step, and simplify your answers. 

  \begin{enumerate}

  \item

    \begin{grading}
      1 point
    \end{grading}

    \begin{problem}[\vfill]
      $\displaystyle 3 e^x+1=5e^x + 6$
    \end{problem}

    \begin{solution}
    \begin{align*}
      3 e^x + 1 &= 5e^x +6 \\
      3  e^x -5e^x &= 6 -1 \\
      -2 e^x &= 5 \\
      e^x &= -\frac{ 5} {2} 
    \end{align*}
      This equation has no solution, as there is no value of $x$ 
      such that $e^x = -\tfrac{5}{2}$.
    \end{solution}

  \item

    \begin{grading}
      3 points. There are multiple ways to solve it; here's a breakdown for the most common method.
      \begin{itemize}[nosep]
      \item 0.5 for taking $\ln$ of both sides
      \item 1 for knowing $\ln \left( 5 \cdot 2^{-x} \right) = \ln 5 + \ln \left( 2^{-x} \right)$
      \item 0.5 for knowing $\ln \left( 2^{-x} \right) = -x \ln 2$
      \item 1 for solving the linear equation $3x = \ln 5 - (\ln 2) x$
      \end{itemize}
    \end{grading}

    \begin{problem}[\stretch{1.5}]
      $e^{3x} = 5 \cdot 2^{-x}$
    \end{problem}

    \begin{align*}
      e^{3x} &= 5\cdot 2^{-x} \\
      \ln (e^{3x}) &= \ln(5\cdot 2^{-x}) \\
      \ln (e^{3x}) &= \ln 5 + \ln(2^{-x}) \\
      3x &= \ln 5 - x\ln 2 \\
      (3+\ln2)x &= \ln 5 \\
      x &= \boxed{\frac{\ln5}{3+\ln2}}
    \end{align*}

  \item

    \begin{grading}
      1 point      
    \end{grading}

    \begin{problem}[\vfill\newpage]
      $(2x + 1)^7 = 3^{14}$
    \end{problem}

    \begin{solution}
      \begin{align*}
        (2x+1)^7 &= 3^{14} \\
        ((2x+1)^7)^{1/7} &= (3^{14})^{1/7} \\
        2x+1 &= 3^2 \\
        2x &= 8 \\
        x &= \boxed{4}
      \end{align*}
    \end{solution}

  \item
    \begin{grading}
      3 points. There are multiple ways to solve it; here's a breakdown for the most common method.
      \begin{itemize}[nosep]
      \item 0.5 for taking $\ln$ of both sides
      \item 0.5 for knowing $\ln \left( \frac{7}{e^{x^2}} \right) = \ln 7 - \ln \left( e^{x^2} \right)$
      \item 0.5 for knowing $\ln \left( e^{x^2} \right) = x^2$
      \item 1.5 for recognizing they have a quadratic equation and correctly using the quadratic formula
      \end{itemize}
    \end{grading}

    \begin{problem}[\vfill]
      $\dfrac{7}{e^{x^2}} = 5^x$
    \end{problem}
    \begin{solution}
      \begin{align*}
        \frac{7}{e^{x^2}} &= 5^x \\
        7 &= 5^x \cdot e^{x^2} \\
        \ln 7 &= \ln\! \left(5^x\cdot e^{x^2}\right) \\
        \ln 7 &= \ln\!\left(5^x\right) + \ln\!\left(e^{x^2}\right) \\
        \ln 7 &= x \ln5 + x^2 \\ 
        0 &= x^2 +(\ln5)x - (\ln 7)
        \intertext{Applying the quadratic formula:}
        x &= \frac{-\ln5 \pm \sqrt{(\ln5)^2-4(1)(-\ln7)}}{2(1)} \\
        x &= \boxed{\frac{-\ln5 \pm \sqrt{(\ln5)^2+4\ln7}}{2}}
      \end{align*}
    \end{solution}

  \item

    \begin{grading}
      3 points. Here's a breakdown for the most common method. 
      \begin{itemize}[nosep]
      \item 1 for recognizing to substitute $u = \ln x$
      \item 1 for correctly solving $u^2 = 2u + 8$
      \item 1 for solving for $x$ once they know $u$
      \end{itemize}
    \end{grading}

    \begin{problem}[\vfill\newpage]
      $(\ln x)^2 = 2 \ln x + 8$
    \end{problem}

    \begin{solution}
      We can substitute $u=\ln x$.
      \begin{align*}
        (\ln x)^2 &= 2\ln x + 8 \\
        u^2 &= 2u + 8 \\
        u^2 - 2u - 8 &= 0 \\
        (u-4)(u+2) &= 0
      \end{align*}
      So $u=\ln x$ is either $4$ or $-2$. This gives us the
      solutions $x=\boxed{e^{4}}$ and $x=\boxed{e^{-2}}$.
    \end{solution}

  \end{enumerate}

\item  

  \begin{enumerate}
  \item

    \begin{grading}
      3 points: 1 for the graph and description of $f(x)$ and 2 for the graph and description of $g(x)$
    \end{grading}

    \begin{problem}[\vfill]
      On the same set of axes, sketch graphs of $f(x) = \log_2(x + 15)$ and $g(x) = 1 + 2 \log_2(-x)$. Explain how each graph is related to the graph of $a(x) = \log_2(x)$. That is, if you start with the graph of $a(x) = \log_2(x)$, what graph transformations could you do to get the graphs of $f$ and $g$? 
    \end{problem}

    \begin{solution}
      The graph of $y=\log_2(x+15)$ is the graph of $y=\log_2(x)$
      shifted to the left by $15$ units. The graph of $y=1+2\log_2(-x)$
      is the graph of $y=\log_2(x)$ reflected over the $y$-axis,
      scaled vertically by a factor of $2$, and shifted up $1$ unit.

      \begin{center}
        \begin{tikzpicture}
          \begin{axis}
            \addplot[domain=-15:5] {ln(x+15)/ln(2)};
            \addplot[domain=-15:0] {1+2*ln(-x)/ln(2)};
          \end{axis}
        \end{tikzpicture}
      \end{center}
    \end{solution}

  \item

    \begin{grading}
      4 points:
      \begin{itemize}[nosep]
      \item 2 points for getting to $x + 15 = 2 (-x)^2$
      \item 1 for solving the quadratic
      \item 1 for recognizing only $x = - \dfrac{5}{2}$ is in the domain. (If they lose this point, their answer is also probably inconsistent with their picture; please point out that inconsistency.) 
      \end{itemize}
    \end{grading}

    \begin{problem}[\vfill\newpage]
      Find all points where the graphs of $f$ and $g$ intersect. (Give the $x$- and $y$-coordinate of each point of intersection.) As usual, please organize your work so that your reader can follow it easily. 
    \end{problem}

    \begin{solution}
      We want to solve the equation $\log_2(x+15) = 1+2\log_2(-x)$.
      We can do this by raising $2$ to both sides.
      \begin{align*}
        \log_2(x+15) &= 1+2\log_2(-x) \\
        2^{\log_2(x+15)} &= 2^{1+2\log_2(-x)} \\
        2^{\log_2(x+15)} &= 2\cdot 2^{2\log_2(-x)} \\
        x+15 &= 2(-x)^2 \\
        x+15 &= 2x^2 \\
        0 &= 2x^2 - x - 15 \\
        0 &= (2x+5)(x-3)
      \end{align*}
      This gives us $x=3$ and $x=-5/2$. However, $x=3$ is not
      in the domain of $g(x)=1+2\log_2(-x)$, so the only
      solution is $\boxed{x=-5/2}$.
    \end{solution}
    

  \end{enumerate}

\end{enumerate}
\psreflection{For next class, read \href{https://people.math.harvard.edu/~jjchen/mathma/docs/handouts/Connally-5-4.html}{\S 5.4 of \emph{Functions Modeling Change} by Connally et. al}.}
\end{document}
